\everymath{\displaystyle}

%%%%%%%%%%%%%%%%%%%%%%%%%%%%%%%%%%%%%%%%
%           main body
%%%%%%%%%%%%%%%%%%%%%%%%%%%%%%%%%%%%%%%%

%-----------------------------------
%	TITLE PAGE
%-----------------------------------

\title[模的直和分解]{\texorpdfstring{\LARGE \bfseries 第9部分\quad 模的直和分解}{第9部分 模的直和分解}}
\author[抽象代数 2]{\Large \bfseries 抽象代数 2}
\date{}

%------------------------------------------------

\begin{document}

%------------------------------------------------

\begin{frame}

\titlepage % Print the title page as the first slide

\end{frame}

%------------------------------------------------

\section[直和]{外直和与内直和}

%------------------------------------------------

\begin{frame}{外直和}

\begin{Def}[外直和]
$M_1, \dots, M_n$ 为 $n$ 个 $R$-模, 令
\[
M_1 \times \dots \times M_n = \{ (\alpha_1, \dots, \alpha_n) \mid \alpha_i \in M_i \}
\]
为 Abel 群, by
\[
(\alpha_1, \dots, \alpha_n) + (\beta_1, \dots, \beta_n) = (\alpha_1 + \beta_1, \dots, \alpha_n + \beta_n).
\]
又 $\forall\ a \in R$, $a(\alpha_1, \dots, \alpha_n) = (a \alpha_1, \dots, a \alpha_n)$,
则 $M_1 \times \dots \times M_n$ 为 $R$-模, 记为 $M_1 \oplus \dots \oplus M_n$,
称为 $M_1, \dots, M_n$ 的外直和.
\end{Def}

\end{frame}

%------------------------------------------------

\begin{frame}{内直和}

\begin{Def}[内直和]
$M$ 为 $R$-模, $N_1, \dots, N_n$ 为 $M$ 的 $n$ 个子模, 若
\[
M = N_1 + \dots + N_n, \qquad
N_i \cap \Big( \sum_{j = 1, j \neq i}^{n} N_j \Big) = 0,
\]
则称 $M$ 为 $N_1, \dots, N_n$ 的内直和, 记为
$M = N_1 \dotplus \dots \dotplus N_n$.
\end{Def}

\end{frame}

%------------------------------------------------

\begin{frame}{可直和分解模}

\begin{Def}[可直和分解模]
若 $M = M_1 \oplus M_2$, 其中 $M_1 \neq 0$, $M_2 \neq 0$,
则称 $M$ 为可直和分解模, 反之则称 $M$ 为不可直和分解模.
\end{Def}

\pause

\begin{eg}
若 $M$ 为 $R$-单模, 则 $M$ 为不可分解的.
\end{eg}

\end{frame}

%------------------------------------------------

\section[幂等元]{嵌入, 投影与幂等元}

%------------------------------------------------

\begin{frame}{嵌入与投影}

若 $M = \bigoplus M_i$, 定义
\[
\delta_i : M_i \to M \quad \text{by } \forall\ \alpha \in M_i,\ \delta_i(\alpha) = \alpha
\qquad (\text{嵌入}),
\]
\[
P_i : M \to M_i \quad \text{by } \forall\ \alpha \in M,\
\alpha = \alpha_1 + \alpha_2 + \dots + \alpha_n\ (\text{唯一}),\
P_i(\alpha) = \alpha_i
\qquad (\text{投影}).
\]

\pause

性质:
\[
P_j \delta_i =
\begin{cases}
0, & \text{若 } i \neq j, \\
1_{M_i}, & \text{若 } i = j.
\end{cases}
\]

\end{frame}

%------------------------------------------------

\begin{frame}{正交幂等元}

$\forall\ 1 \leq i \leq n$, 令 $e_i = \delta_i P_i \in \mathrm{End}_R(M)$,
\[
e_i|_{M_i} = 1_{M_i}, \qquad e_i|_{M_j} = 0 \ (j \neq i),
\]
$\{ e_1, e_2, \dots, e_n \} \subset \mathrm{End}_R(M)$, $e_i \neq 1$, $e_i \neq 0$:
\begin{itemize}
\item \circled{1} $e_i^2 = e_i$, $1 \leq i \leq n$;
\item \circled{2} $e_i e_j = 0$, 若 $i \neq j$;
\item \circled{3} $\sum_i e_i = 1_M$.
\end{itemize}

\end{frame}

%------------------------------------------------

\begin{frame}{非平凡的正交幂等元}

结论: 若 $M$ 可直和分解, $M = \bigoplus M_i$, 则 $\mathrm{End}_R(M)$ 中存在一组非平凡幂等元
$\{ e_1, \dots, e_n \}$ 满足 $e_i e_j = 0\ (i \neq j)$, $\sum e_i = 1$,
称为 $\mathrm{End}_R(M)$ 中一组非平凡的正交幂等元.

\end{frame}

%------------------------------------------------

\begin{frame}{不可分解的刻画}

\begin{lem}[不可分解的刻画]
$M$ 为一个 $R$-模, $M$ 为不可直和分解的 $R$-模 $\iff$
$\mathrm{End}_R(M)$ 仅有平凡幂等元.
\end{lem}

% 注 (转写): 手稿证明只写了单向 (存在非平凡幂等元则可分解), 见升级清单.

\startproof[bg=yellow, fg=red](证明)
若 $\mathrm{End}_R(M)$ 中有 $e \in \mathrm{End}_R(M)$ 满足 $e^2 = e$, $e \neq 1$, $e \neq 0$,
那么 $(1 - e)^2 = 1 - e$, $1 - e \neq 1$, $1 - e \neq 0$,
即 $e$, $1 - e$ 都是 $\mathrm{End}_R(M)$ 中的非平凡幂等元.
令 $M_1 = e(M)$, $M_2 = (1 - e)(M)$,
易验证知 $M = M_1 \oplus M_2$.
\qed

\end{frame}

%------------------------------------------------

\section[局部环]{局部环与强不可分解}

%------------------------------------------------

\begin{frame}{局部环与强不可分解}

\begin{Def}[(非交换)局部环]
$R$ 为幺环, $U(R)$ 为单位群, 若 $I = R \setminus U(R)$ 为 $R$ 的一个理想,
则称 $R$ 为一个(非交换)局部环.
\end{Def}

\pause

\begin{Def}[强不可分解模]
$M$ 为 $R$-模, $\mathrm{End}_R(M)$ 为一个局部环,
则称 $M$ 为一个强不可分解的 $R$-模.
\end{Def}

\end{frame}

%------------------------------------------------

\begin{frame}{强不可分解的唯一性}

\begin{lem}[强不可分解的唯一性]
$R$-模 $M$ 可以分解为有限个强不可分解子模的直和, 那么 $M$ 可以分解为有限多个
不可分解子模的直和, 且分解唯一.
\end{lem}

\startproof[bg=yellow, fg=red](证明)
若 $M = M_1 \oplus M_2 \oplus \dots \oplus M_n$, 其中 $M_i$ 为强不可分解的子模.
要证明对 $M$ 的任一分解 $M = N_1 \oplus \dots \oplus N_m$, 其中 $N_i$ 为不可分解模,
那么 $m = n$, 且 $\forall\ 1 \leq i \leq m$,
$N_i \cong M_j$ (某个 $j$, $1 \leq j \leq n$).
\qed

\end{frame}

%------------------------------------------------

\begin{frame}{强不可分解的唯一性: 归纳计划与断言}

\startproof[bg=yellow, fg=red](证明, 续)
$*$ 对 $n$ 作归纳证明, 再加上一个断言:
\[
\text{若 } M \xrightarrow{\ f\ } N \xrightarrow{\ g\ } M \text{ 为 } R\text{-模同态},\
N \text{ 为强不可分解的,}
\]
\[
\text{若 } g \circ f \in \mathrm{End}_R(M) \text{ 为同构, 则 } g \text{ 与 } f \text{ 都为同构.}
\]
\qed

\pause

\begin{itemize}
\item 结论:
\item \circled{1} 有限长度模, 强不可分解 $\Longleftarrow$ 不可分解;
\item \circled{2} 有限长度模 Krull--Schmidt 分解定理.
\end{itemize}

\end{frame}

%------------------------------------------------

\section[Fitting]{Fitting 直和分解}

%------------------------------------------------

\begin{frame}{Fitting 直和分解}

模的具体的直和分解 (Fitting 直和分解):
若 $M$ 为 $R$-模, 任一 $\sigma \in \mathrm{End}_R(M)$, 令 $\sigma^0 = 1_M$, 则
\[
M = \sigma^0(M) \supset \sigma(M) \supset \dots \supset \sigma^n(M) \supset \dots
\]
\[
0 = \ker(\sigma^0) \subset \ker(\sigma) \subset \dots \subset \ker(\sigma^n) \subset \dots
\]
令
\[
\sigma^{\infty}(M) = \bigcap_{n=0}^{\infty} \sigma^n(M), \qquad
\sigma^{-\infty}(M) = \bigcup_{n=0}^{\infty} \ker(\sigma^n).
\]

\end{frame}

%------------------------------------------------

\begin{frame}{Fitting 直和分解定理}

\begin{lem}[Fitting 直和分解]
$M$ 为有限长度模, 则 $\forall\ \sigma \in \mathrm{End}_R(M)$,
\[
M = \sigma^{\infty}(M) \oplus \sigma^{-\infty}(M),
\]
且 $\sigma|_{\sigma^{\infty}(M)}$ 为 $\sigma^{\infty}(M)$ 的同构,
$\sigma|_{\sigma^{-\infty}(M)}$ 为幂零元素.
\end{lem}

\end{frame}

%------------------------------------------------

\begin{frame}{Fitting 直和分解的证明 (一)}

\startproof[bg=yellow, fg=red](证明)
$M$ 为有限长度模, 则
$\exists\ k \geq 1$, 使得 $\forall\ n \geq k$, $\sigma^n(M) = \sigma^k(M)$;
$\exists\ t \geq 1$, 使得 $\forall\ n \geq t$, $\ker(\sigma^n) = \ker(\sigma^t)$.
令 $n = \max\{ k, t \}$, 有 $\sigma^{\infty}(M) = \sigma^n(M)$,
$\sigma^{-\infty}(M) = \ker(\sigma^n)$.
仅需证明 $M = \sigma^n(M) \oplus \ker(\sigma^n)$, 即:
\circled{1} $M = \sigma^n(M) + \ker(\sigma^n)$;
\circled{2} $\sigma^n(M) \cap \ker(\sigma^n) = 0$.

\pause

对\circled{2}: $\forall\ \alpha \in \sigma^n(M) \cap \ker(\sigma^n)$,
$\exists\ x \in M$, 使得 $\alpha = \sigma^n(x)$,
那么 $0 = \sigma^n(\alpha) = \sigma^{2n}(x) = \sigma^n(x) = \alpha$.
\qed

\end{frame}

%------------------------------------------------

\begin{frame}{Fitting 直和分解的证明 (二)}

\startproof[bg=yellow, fg=red](\circled{1} 的证明)
$\forall\ \alpha \in M$, 令 $\alpha_1 = \sigma^n(\alpha)$,
$\alpha_1 \in \sigma^n(M) = \sigma^{2n}(M)$,
那么 $\exists\ x \in M$, 使得 $\alpha_1 = \sigma^{2n}(x)$,
即 $\sigma^n(\alpha) - \sigma^{2n}(x) = 0$,
即 $\sigma^n(\alpha - \sigma^n(x)) = 0$,
那么 $\alpha - \sigma^n(x) \in \ker(\sigma^n)$,
令 $\alpha_2 = \alpha - \sigma^n(x)$, 即有 $\alpha = \alpha_1 + \alpha_2$.
\qed

\end{frame}

%------------------------------------------------

\begin{frame}{Fitting 直和分解的证明 (三)}

\startproof[bg=yellow, fg=red](2° 与 3° 的证明)
$2^{\circ}$. $\sigma(\sigma^n(M)) = \sigma^{n+1}(M) = \sigma^n(M)$,
即知 $\sigma$ 在 $\sigma^n(M)$ 为满的.
又任取 $\alpha, \beta \in \sigma^n(M)$, 若 $\sigma(\alpha) = \sigma(\beta)$,
那么 $\sigma(\alpha - \beta) = 0$,
$\alpha - \beta \in \ker \sigma \subset \ker \sigma^n$,
$\because$ $\ker \sigma^n \cap \sigma^n(M) = 0$, 所以 $\alpha = \beta$,
即 $\sigma$ 在 $\sigma^n(M)$ 为单的,
即知 $\sigma|_{\sigma^n(M)} = \sigma|_{\sigma^{\infty}(M)}$ 为同构.

\pause

$3^{\circ}$. $\forall\ \alpha \in \ker(\sigma^n)$, $\sigma^n(\alpha) = 0$,
即知 $\sigma|_{\ker \sigma^n}$ 为幂零元素,
即 $\sigma|_{\sigma^{-\infty}(M)}$ 为幂零元素.
\qed

\end{frame}

%------------------------------------------------

\section[分解定理]{Krull--Schmidt 分解定理}

%------------------------------------------------

\begin{frame}{强不可分解与不可分解}

\begin{cor}[强不可分解与不可分解]
$M$ 为一个有限长度模, 那么 $M$ 强不可分解 $\iff$ $M$ 不可解.
\end{cor}

% 注 (转写): 手稿原文「不可解」, 应为「不可分解」, 见升级清单.
% 注 (转写): 手稿原文 U(End_R(R)), 括号内应为 M, 见升级清单.

\startproof[bg=yellow, fg=red](证明)
$M$ 为不可分解的, 我们证明 $\mathrm{End}_R(M)$ 为 Local 环,
仅需证 $\mathrm{End}_R(M) \setminus U(\mathrm{End}_R(R)) = I$ 为 $\mathrm{End}_R(M)$ 的理想.
因为 $\forall\ \tau \in \mathrm{End}_R(M)$, 有 Fitting 分解
$M = \tau^n(M) \oplus \ker(\tau^n)$,
$M$ 不可分解, 那么有 $\tau^n(M) = M$ 或 $\ker(\tau^n) = M$.
若 $\tau^n(M) = M$, 则 $\tau$ 为可逆元; 若 $\ker(\tau^n) = M$, 则 $\tau$ 为幂零元素.
所以 $\mathrm{End}_R(M)$ 中元素为同构或者幂零的.

\end{frame}

%------------------------------------------------

\begin{frame}{强不可分解与不可分解 (续)}

\startproof[bg=yellow, fg=red](证明, 续)
任取 $\alpha \in I$, $\alpha$ 为幂零元素, 那么 $\forall\ \tau \in \mathrm{End}_R(M)$,
$\alpha \tau$, $\tau \alpha$ 不是同构, 因此 $\alpha \tau$, $\tau \alpha \in I$.
任取 $\alpha, \beta \in I$, 若 $\alpha + \beta \notin I$, 那么 $\alpha + \beta$ 可逆,
令 $\alpha_1 = \alpha \cdot (\alpha + \beta)^{-1}$,
$\alpha_2 = \beta (\alpha + \beta)^{-1}$, 有 $\alpha_1 + \alpha_2 = 1$.
由于 $\alpha_1$ 与 $\alpha_2 \in I$, 为幂零元素, 令 $\alpha_2^n = 0$,
那么 $(1 - \alpha_2)(1 + \alpha_2 + \dots + \alpha_2^{n-1}) = 1 - \alpha_2^n = 1$,
那么 $1 - \alpha_2 = \alpha_1 \in U(\mathrm{End}_R(M))$, 这与 $\alpha_1 \in I$ 矛盾.
此即可知 $I$ 为理想.
\qed

\end{frame}

%------------------------------------------------

\begin{frame}{Krull--Schmidt 定理}

% 注 (转写): 手稿把 Krull 拼作 Knell (两处), slides 按人名标准拼写, 见升级清单.

\begin{thm}[Krull--Schmidt 定理]
任一有限长度的 $R$-模 $M$ 可以分解为有限个不可分解模的直和, 而且分解是唯一的.
\end{thm}

\end{frame}

%------------------------------------------------

\begin{frame}{Krull--Schmidt 定理的证明}

\startproof[bg=yellow, fg=red](证明)
$M$ 为有限长度的 $R$-模, 长度为 $\ell(M)$.
$N \subset M$ 为子模, 也是有限长度的, $\ell(N) \leq \ell(M)$.
对 $\ell(M)$ 作归纳, $n = 1 \iff M$ 为单模, 命题为真.
假若 $\ell(M) \leq n - 1$ 时命题为真, 那么当 $\ell(M) = n$ 时:

\pause

若 $M$ 可以分解, 令 $M = M_1 \oplus M_2$,
那么 $\ell(M_1) < \ell(M)$, $\ell(M_2) < \ell(M)$,
由归纳假设,
\[
M_1 = N_1 \oplus N_2 \oplus \dots \oplus N_s, \qquad
M_2 = H_1 \oplus H_2 \oplus \dots \oplus H_t,
\]
那么 $M = N_1 \oplus \dots \oplus N_s \oplus H_1 \oplus \dots \oplus H_t$,
且 $N_i$, $H_j$ 为强不可分解的,
由强不可分解的唯一性引理即知唯一性.
若 $M$ 不可分解, 则 $M$ 为强不可分解的,
由强不可分解的唯一性引理知命题为真.
\qed

\end{frame}

%------------------------------------------------

\end{document}
